% Seção: Funcionalidades e resultados
% Dono: cada dono de feature · Orçamento: 4 páginas no total · plano-de-projeto §14.1
% Conteúdo: uma subseção por feature, com tela, resultado e complexidade.
%
% Como encaixar a sua feature: crie relatorio/secoes/<ID>.tex começando com
% \subsection{<título>}\label{sec:<ID>}. Ela entra aqui sozinha, sem editar este arquivo.
% Essenciais mostram um aviso enquanto o arquivo não existir; desejáveis e opcionais
% só aparecem quando o arquivo existir.
% Cada subseção traz: figura (tela), resultado com o número medido, complexidade e
% uma limitação. Nenhum número sem execução real (AGENTS.md §5, regra 6).
\section{Funcionalidades e resultados}\label{sec:funcionalidades}

% Essenciais (P1)
\secaofeature{F-ROTA}{Trajeto mais rápido}
\secaofeature{F-ESPERA}{Espera e simulação de trajeto}
\secaofeature{F-BALD}{Menos baldeações}
\secaofeature{F-MET}{Centralidades}
\secaofeature{F-EFI}{Eficiência da rede}
\secaofeature{F-RES}{Falhas e resiliência}

% Desejáveis (P2)
\secaofeatureopcional{F-ROTA-2}
\secaofeatureopcional{F-ESPERA-2}
\secaofeatureopcional{F-BALD-2}
\secaofeatureopcional{F-MET-2}
\secaofeatureopcional{F-EFI-2}
\secaofeatureopcional{F-RES-2}

% Opcionais
\secaofeatureopcional{F-ROTA-OPC}
\secaofeatureopcional{F-ESPERA-OPC}
\secaofeatureopcional{F-BALD-OPC}
\secaofeatureopcional{F-MET-OPC}
\secaofeatureopcional{F-EFI-OPC}
\secaofeatureopcional{F-RES-OPC}
